\documentclass[11pt]{article}

\usepackage[margin=1in]{geometry}
\usepackage{amsmath}
\usepackage{amssymb}
\usepackage{booktabs}
\usepackage{array}
\usepackage{graphicx}
\usepackage{xcolor}
\usepackage{caption}
\usepackage{microtype}
\usepackage[colorlinks=true,linkcolor=blue!45!black,citecolor=blue!45!black,urlcolor=blue!45!black]{hyperref}

\newtheorem{theorem}{Theorem}
\newenvironment{proofenv}{\noindent\textit{Proof.}\ }{\hfill$\square$\par\medskip}

\newcommand{\code}[1]{\texttt{#1}}
\newcommand{\usff}{\code{us54}}

\title{FishAI v0.5: Measuring Play Style Without Skill\\ in a 54-Card Literature Variant}
\author{Allen Hsieh\\[2pt] {\small FishAI project}}
\date{August 2026}

\begin{document}
\maketitle

\begin{abstract}
\noindent
FishAI v0.5 is a style-conditioned engine and simulation laboratory for Canadian Fish
(Literature) under the \usff{} rule dialect: 54 cards, nine half-suits of six, out-of-turn
declares, any declaration error awarded to the opponents, and a race to five sets. Nine
parameterized play styles share one identical full-strength inference engine, so that
differences in outcome measure \emph{style} rather than \emph{skill}. The roster is evaluated
in a full-precision duplicate-deal round-robin --- 36 cells, 4{,}300 mirrored deal pairs per
cell, 309{,}600 games, per-cell standard error at most $0.005$ --- and the payoff matrix is
analyzed with mean score, maximin, Bradley--Terry, Nash averaging, $\alpha$-Rank, and a Hodge
decomposition into transitive and cyclic parts. The matrix is almost purely transitive (cyclic
energy $0.0112$ against a $0.15$ threshold), and one style, \emph{Punter}, satisfies all four
pre-stated dominance criteria: highest mean score ($0.5829$), maximin above one half
($0.5190$), negligible cyclic energy, and exploitability $0.0025$ against a rivals' median of
$0.0444$. Two measured caveats qualify the headline and are reported as first-class results.
First, the declare-threshold axis that the roster's aggressive-to-conservative labels
advertise is \emph{inert} across the range the roster spans: the shared engine's confidence
estimates are bimodal, so every threshold in the occupied band selects an identical set of
declares, and Punter wins through its ask-scoring gamble bonus and uncertainty tolerance, not
through a lower declare bar. Second, the card-hoarding style's signature benefit --- a
repeatable turn-pass through a team-contained book --- was derived, implemented, and measured
to be worth nothing: the Hoarder finishes 6th of 9 with and without it, and the best measured
counter to hoarding is simply not hoarding. To the best of our knowledge, resting on an
internal literature synthesis that found no peer-reviewed prior work on the game, this is the
first quantitative game-AI study of Literature.
\end{abstract}

\section{Introduction}
\label{sec:intro}

Literature (Canadian Fish) is a six-player, two-team card game of public asks and hidden
hands, in which every action is broadcast to the whole table and inference --- not chance ---
decides who knows what. Communities that play it argue, persistently and without data, about
temperament: whether the aggressive declarer beats the certainty purist, whether hoarding
ask-licences is wisdom or superstition, whether the styles simply counter one another in a
rock--paper--scissors loop. FishAI v0.5 was built to make that argument quantitative for one
concretely pinned dialect. The research question, fixed before any tournament ran, is:

\begin{quote}
\emph{Is there a quantitatively superior play style --- or do styles counter each other?}
\end{quote}

The question is harder to pose honestly than it looks, for two reasons. The first is a
confound: a ``style'' implemented as a weaker bot measures nothing but its author's effort.
FishAI's design rule is therefore that every style shares one identical, full-strength
inference engine and differs only in the policy layer above it (Section~\ref{sec:style}).
The second is that ``best'' is a meta-game notion, and meta-games are often intransitive:
ranking systems that assume transitive skill, such as Elo, are actively misleading when the
payoff matrix contains cycles \cite{balduzzi2018}. The laboratory therefore decomposes the
matrix into transitive and cyclic parts and commits, in advance, to a four-criterion decision
rule under which ``no dominant style; here is the counter-structure'' is a first-class
outcome, not a failure (Section~\ref{sec:method}).

This paper documents the v0.5 system and its measurements. Its contributions are:

\begin{enumerate}
\item \textbf{A rule-complete specification and engine for the \usff{} dialect} --- the
  54-card ``US student'' variant --- with its derived consequences: the sign flip in
  declaration risk, a tie-impossibility theorem, and the observation that whole-game
  termination is a property of the acting policy rather than of the rules
  (Section~\ref{sec:game}).
\item \textbf{A style-not-skill experimental design}: nine parameterized styles over one
  shared inference engine, a decision-divergence instrument that verifies each style is
  measurably distinct from the control, and duplicate-deal evaluation with shared seed sets
  and seat rotation (Sections~\ref{sec:style}--\ref{sec:method}).
\item \textbf{A full-precision tournament and its verdict}: 309{,}600 games over 36 cells,
  analyzed side by side with mean score, maximin, Bradley--Terry, Nash averaging,
  $\alpha$-Rank and Hodge decomposition, yielding a \emph{dominant} verdict for the Punter
  style under all four pre-stated criteria, together with the per-style diagnostics that
  explain the win (Section~\ref{sec:results}).
\item \textbf{Two negative results, reported as results.} The declare-threshold axis is inert
  across the roster's range, with a measured bimodal-confidence mechanism
  (Section~\ref{sec:inert}); and the Hoarder style's thesis is refuted end-to-end --- its
  costs measured, its theorized benefit implemented, and the benefit measured to be worth
  nothing (Section~\ref{sec:hoarder}).
\end{enumerate}

Throughout, every quantitative claim is traceable either to the committed tournament artifact
\cite{artifact} or to the repository's measurement documents
\cite{rules,styles,botlab,containment}; Section~\ref{sec:provenance} records the artifact's
provenance digests.

\section{Related work}
\label{sec:related}

\paragraph{Literature as a game-AI domain.}
An internal research synthesis for this project \cite{playstyles} --- a 65-style taxonomy of
Literature play compiled with explicit evidence tiers --- reports that targeted searches
across arXiv, the ACM Digital Library, IEEE Xplore, AAAI proceedings, ResearchGate and
university thesis repositories found \emph{zero peer-reviewed papers, theses or technical
reports on Literature as a game-AI domain}: no baseline agent, no published environment, no
state-space or equilibrium analysis, no benchmark. That synthesis's own access to non-GitHub
sources was constrained and it flags many of its secondary citations as unverified; we treat
its absence claim as a careful search's negative report, current as of mid-2026, rather than
as proof. The only public Literature reinforcement-learning project we are aware of is a
four-player Q-learning implementation that reports no win rates \cite{somani}. The deepest
strategy text on the game remains Develin's chapter \cite{develin}, which independently
attests the contained-book result used in Section~\ref{sec:hoarder}.

\paragraph{Evaluating populations of agents.}
The evaluation stack is assembled from the multi-agent evaluation literature: Nash averaging,
which plays a meta-game on the evaluation data itself and is robust to redundant entrants
\cite{balduzzi2018}; $\alpha$-Rank, an evolutionary ranking that survives intransitivity
\cite{omidshafiei2019}; and the Hodge-theoretic decomposition of a pairwise-comparison matrix
into gradient (transitive) and cyclic parts \cite{jiang2011,balduzzi2018}. Bradley--Terry
ratings \cite{bradleyterry1952} are reported alongside as the classical transitive fit.
Duplicate (mirrored) deals are the card-game form of common random numbers, standard in
duplicate bridge; the poker literature's stronger variance-reduction machinery
\cite{burch2018} was measured to be unnecessary at this game's deal-dependence
(Section~\ref{sec:duplicate}). Sequential probability ratio tests \cite{wald1945} gate the
tuning and exploitability searches, following modern chess-engine testing practice; the
published matrix itself is fixed-$N$, since sequential stopping biases effect-size estimates.
False-discovery control over the 36 simultaneous cells uses Benjamini--Hochberg
\cite{benjamini1995}.

\paragraph{Imperfect information and team play.}
Formally, Literature is a two-team constant-sum imperfect-information game whose teams may
coordinate only ex ante --- the regime whose solution concept is team-maxmin equilibrium with
correlation \cite{celli2018}. All hidden information originates at the single chance node of
the deal, and every subsequent event is public; the synthesis \cite{playstyles} notes that
branching is tiny (at most 120 legal asks, typically 20--40) while one seat's deal-time
information set in the 48-card baseline is on the order of $10^{24}$, so belief fidelity, not
search depth, is the game's difficulty. Search-based agents --- ISMCTS \cite{cowling2012} and
perfect-information Monte Carlo \cite{long2010} --- are deliberately out of scope for v0.5: a
search bot has a budget, not a style, and it would answer a different question. PIMC in
particular must never evaluate declarations here, since strategy fusion makes every sampled
world report the declaration succeeding \cite{long2010,containment}. The style-vs-style
setting also inherits the coordination literature's central warning: conventions formed in
self-play need not transfer to novel partners \cite{hu2020,bard2020}, which is why cross-play
against an independently written bot is part of the design --- and, not yet having been run,
is listed as an open threat in Section~\ref{sec:threats}. The shared inference engine itself
belongs to the constraint-propagation family developed for deduction games such as Clue
\cite{neller2018}.

\section{The game and the \usff{} dialect}
\label{sec:game}

\subsection{Base game}

Six players in two teams of three, seated alternately, play with a deck partitioned into
six-card \emph{half-suits} (``sets'' or ``books''). On your turn you ask one specific card of
one specific opponent; you must yourself hold a card of the asked set, and may not ask for a
card you hold. A hit transfers the card face up and you ask again; a miss passes the turn to
the player you asked. At any point a team may \emph{declare} a set by naming the exact holder
of all six cards, every assignment being a seat on the declarer's own team. Card counts and
the full log of asks, results and declares are public. The team with more sets wins.

\subsection{The \usff{} dialect}

FishAI v0.5 plays the ``US student'' dialect, specified in full with derived consequences and
test vectors in the repository's rules document \cite{rules}. Table~\ref{tab:dialect}
summarizes the deltas against the 48-card baseline of the main reference rules text
\cite{pagat}.

\begin{table}[t]
\centering
\caption{The \usff{} dialect against the 48-card baseline \cite{rules}. The 48-card
wrong-declare outcome splits by error type (its rows 14--15): a set any opponent holds a card
of is scored by the opponents --- the same outcome as under \usff{} --- and only a set the
declarer's own team wholly holds, misassigned, is void.}
\label{tab:dialect}
\small
\begin{tabular}{@{}lll@{}}
\toprule
 & \usff{} (this work) & 48-card baseline \\
\midrule
Deck & 54 --- standard 52 plus two jokers & 48 (8s removed) \\
Sets & 9 half-suits of 6; the ninth is
  $8\clubsuit\,8\diamondsuit\,8\heartsuit\,8\spadesuit$ + both jokers & 8 \\
Hand & 9 cards & 8 cards \\
Declaring & any time, out of turn, in a declare window & own turn only \\
Wrong declare & \textbf{the opponents score the set} &
  \begin{tabular}[t]{@{}l@{}}opponent-held: they score;\\ misassigned: void\end{tabular} \\
Game end & the moment a team is awarded 5 sets & when all sets are resolved \\
Ties & impossible (Theorem~\ref{thm:tie}) & possible (4--4) \\
\bottomrule
\end{tabular}
\end{table}

Out-of-turn declaring is resolved deterministically through a \emph{declare window}: after
every action, each seat in cyclic order from the turn-holder is offered the chance to declare;
declining is itself a move, and when the turn-holder has no legal ask, declining is illegal
(\code{MUST\_DECLARE}) --- a rule that guarantees a declare window can never stall
\cite{rules}. Two of the dialect's rows change the strategy space far more than they appear
to, and the next two subsections make them precise.

\subsection{The sign flip in declaration risk}
\label{sec:signflip}

Under the 48-card baseline, misassigning a set your own team holds produces a \emph{void}:
zero for everyone. Under \usff{} any error at all awards the set to the opponents. Relative to
simply holding the set, a failed declaration used to cost one point (yours, burned); it now
costs one point \emph{and} hands the opponents one --- a two-point swing in a race to five
\cite{styles}. Writing $c$ for the probability the set is contained by your team and $a$ for
the probability your seat-assignment is exactly right given containment, the expected-value
thresholds differ as in Table~\ref{tab:ev} \cite{containment}.

\begin{table}[t]
\centering
\caption{Declaration expected-value thresholds under the two rule sets \cite{containment}.}
\label{tab:ev}
\small
\begin{tabular}{@{}lll@{}}
\toprule
Rule set & Outcome of any error & Declare iff \\
\midrule
48-card baseline & misassigned: void; opponent-held: they score & $c > 1/(1+a)$ \\
\usff{} & opponents score ($-1$ us, $+1$ them) & $c \cdot a > 1/2$ \\
\bottomrule
\end{tabular}
\end{table}

Every declaration threshold tuned against the 48-card game is therefore too loose under
\usff{}. The roster's thresholds are treated as 48-card \emph{appetites} and mapped through
$t_{\mathrm{us54}} = (1 + t_{\mathrm{pagat48}})/2$ before shipping --- Blitz $0.70 \to 0.85$,
Punter $0.55 \to 0.775$, Hoarder $0.95 \to 0.975$ \cite{styles}. Section~\ref{sec:inert}
reports the measured, and unexpected, fate of this entire axis.

\subsection{Ties are impossible --- and termination is a policy property}

\begin{theorem}[Tie impossibility]
\label{thm:tie}
Under \usff{}, no final state is level: every game that terminates ends in a clinch at
exactly five sets, and a tie can never be emitted.
\end{theorem}

\begin{proofenv}
Because the void outcome is abolished, every resolved set is awarded to exactly one team.
There are nine sets. If neither team has reached five, each holds at most four, so at most
eight sets are resolved --- hence if all nine were resolved, some team reached five first, a
contradiction with the game still running. The clinch target $\lfloor 9/2 \rfloor + 1 = 5$ is
therefore always reached before the sets run out, and no final state can be level
\cite{rules}.
\end{proofenv}

The qualifier in the theorem is deliberate. The pigeonhole argument guarantees that declares
cannot run out before a clinch; it does not force declares to happen. Literature can
\emph{livelock}: an adversarial ask policy that always takes the last legal ask ping-pongs
misses between two seats forever, with the position repeating exactly --- no rule violated,
every state holding a legal action. This reproduces under both rule sets and is a property of
the game, not the dialect \cite{rules}. The architectural consequence is that whole-game
termination belongs to the acting policy: the simulator's step cap (5{,}000) is load-bearing
and instrumented, and in the reported tournament it never fired
(Section~\ref{sec:provenance}). The engine itself is pure, deterministic TypeScript --- the
same seed reproduces a byte-identical game --- and adversarially audited: a 250-game lockstep
differential against the pre-variant engine (8{,}155 reduce steps, zero state mismatches),
130{,}480 deliberately illegal actions confirming stable error codes, and a 10{,}000-game
\usff{} fuzz gate asserting that every game terminates in a clinch at exactly five sets with
no tie ever emitted \cite{repo}.

\section{Measuring style without skill}
\label{sec:style}

\subsection{The design rule}

The central confound in any style comparison is that a ``style'' may simply be a weaker bot.
FishAI's rule: \textbf{every style shares one identical, full-strength inference engine} ---
a constraint-propagation knowledge layer over the public log --- and styles differ only in the
parameterized policy above it \cite{botlab}. Bots consume only \code{SeatView}, the exact
public-information payload a human in that seat would see, and a test enforces that no other
hand is reachable: cheating is structurally impossible rather than policed \cite{repo}. One
guardrail is imposed on every style: the certainty bonus in ask scoring stays at or above its
baseline, because below it a style can rank an uncertain ask above a certain hit --- which is
not a style but a bug that would dominate the results \cite{styles}.

\subsection{The parameter surface}

A style is a \code{StyleParams} vector \cite{styles} over six groups of knobs:

\begin{itemize}
\item \emph{Ask scoring}: weights on hit probability (\code{wHit}), progress toward
  nearly-secured sets (\code{wProgress}), information gain (\code{wNarrow}); a floor on
  acceptable hit probability (\code{minHitP}); a bonus for asks that would complete a set
  (\code{gambleBonus}); and the protected certainty bonus.
\item \emph{Declare policy}: the confidence bar (\code{declareThreshold}), tolerated guessed
  cards (\code{declareMaxUncertain}), and the certainty-only and own-hand-only extremes.
\item \emph{Declare window} (new under \usff{}): how early in the window to fire
  (\code{declareEagerness}); whether and at what separate bar to declare sets the seat holds
  no card of (\code{foreignDeclare}, \code{foreignDeclareThreshold}).
\item \emph{Clinch} (new under \usff{}): preference for the closing declare at four sets
  (\code{clinchAggression}) and weight on denying the opponents' fifth (\code{denialWeight}).
\item \emph{Information and targeting}: the leak-protection tiebreak (\code{leakEpsilon},
  \code{leakThreshold}), signalling asks, and whom to promote on a likely miss
  (\code{missTarget}).
\item \emph{Hoarding and the turn-pass}: keep ask-licences in at least $N$ sets
  (\code{hoardBooks}), refuse declares that shrink the hand below $N$ (\code{minHandSize}),
  and the contained-book turn-pass appetite (\code{containedPass},
  Section~\ref{sec:containedpass}).
\end{itemize}

\subsection{The roster}
\label{sec:roster}

Nine styles (Table~\ref{tab:roster}). Their full shipped vectors are embedded in the committed
artifact \cite{artifact}; the table lists each style's defining deviations from the Balanced
control.

\begin{table}[t]
\centering
\caption{The nine-style roster. Settings are the shipped \usff{} values; declaration
thresholds are 48-card appetites mapped through $(1+t)/2$ (Section~\ref{sec:signflip}).}
\label{tab:roster}
\footnotesize
\begin{tabular}{@{}p{1.5cm}p{1.7cm}p{4.6cm}p{6.6cm}@{}}
\toprule
Style & Family & Thesis & Defining settings (shipped) \\
\midrule
Balanced & control & the tuned \usff{} baseline every other style is read against &
  baseline weights; \code{declareThreshold} 0.90; \code{containedPass} 1 \\
Blitz & aggressive & tempo and sets now; information is cheap &
  \code{wHit} 90, \code{wProgress} 30, \code{declareThreshold} 0.85,
  \code{declareEagerness} 0.9, \code{leakEpsilon} 0, no signalling,
  \code{missTarget} \code{most} \\
Punter & aggressive & chase the completing card; accept the gift risk &
  \code{gambleBonus} 25, \code{minHitP} 0, \code{declareThreshold} 0.775,
  \code{declareMaxUncertain} 2 \\
Banker & conservative & never gift a set; bank only certainties &
  \code{declareOnlyWhenCertain}, \code{minHitP} 0.25, \code{declareEagerness} 0.2 \\
Turtle & passive & minimum risk; declare only sets wholly in hand &
  \code{declareOnlyOwnHand}, \code{minHitP} 0.4, no signalling, no foreign declares \\
Hoarder & optionality & keep ask-licences, stay alive, delay &
  \code{hoardBooks} 3, \code{minHandSize} 2, \code{declareThreshold} 0.975,
  \code{declareEagerness} 0.1, \code{containedPass} 1.33 \\
Scout & information & deduce first, collect later &
  \code{wNarrow} 40, \code{wHit} 55, \code{declareOnlyWhenCertain} \\
Ghost & information & deny opponents the read &
  \code{leakEpsilon} 6, \code{leakThreshold} 3, no signalling \\
Archivist & information & track sets you hold nothing in; declare them for teammates &
  \code{foreignDeclare} at threshold 0.95, \code{wNarrow} 30,
  \code{declareEagerness} 0.7 \\
\bottomrule
\end{tabular}
\end{table}

Two roster notes matter for reading everything that follows. First, the family labels narrate
an aggressive-to-conservative spectrum of \emph{declaring}; Section~\ref{sec:inert} shows by
measurement that the axis those labels advertise does not fire anywhere in the roster's range,
so the labels must be read as naming the styles' \emph{appetites}, not their realized
declaring behaviour. The styles are nonetheless measurably distinct
(Section~\ref{sec:distinct}) --- through their ask policies, their gating extremes, and their
window and hoarding knobs. Second, the Archivist is the style \usff{} created: out-of-turn
declaring makes it worthwhile to spend inference on sets one can never ask into, purely to
declare them for teammates.

\section{Methodology}
\label{sec:method}

\subsection{Duplicate deals, shared seeds, seat rotation}
\label{sec:duplicate}

One trial is one seeded deal played in both orientations --- the style assignment swapped
between the teams, the start seat held identical --- so that a lucky deal cancels within the
pair, exactly as in duplicate bridge; this is the common-random-numbers estimator. Measured on
this engine at design time (500 mirrored pairs), pairing reduced the variance of the score
estimate by a factor of $1.36$ relative to independent games \cite{botlab}; the realized
per-cell variance ratios in the reported tournament average $1.33$ (range $1.14$--$2.23$)
\cite{artifact}. This is far below the $30\times$ reductions achievable in poker with
AIVAT-style control variates \cite{burch2018} --- expected, because the whole deck is dealt
and outcomes are far less deal-dominated --- and the measured gain is why no heavier machinery
was built. Every cell uses the same seed list (prefix \code{style-v1}), so cross-cell
comparisons share deals too; start seats rotate $i \bmod 6$; and because the policies are
deterministic, sample diversity comes entirely from seeds --- the harness asserts
$\code{distinctSeeds} = \text{pairs}$.

\subsection{Precision and primary metric}

The published matrix is fixed-$N$ at 4{,}300 pairs (8{,}600 games) per cell --- above the
2{,}600-pair budget derived for a target standard error of $0.005$ --- and the largest
realized cell SE is exactly $0.0050$ \cite{artifact}. The primary metric is score rate (win
$=1$, tie $=\tfrac12$, loss $=0$), adopted because the 48-card game produces ties in roughly a
quarter of mirror games \cite{botlab}. Under \usff{} ties are impossible
(Theorem~\ref{thm:tie}), so score rate coincides with win rate; the artifact asserts zero ties
across all 309{,}600 games rather than rendering a column that cannot populate. Sequential
tests are used only where they belong: SPRT ($\alpha=\beta=0.05$) gates tuning and
exploitability-search candidates \cite{wald1945}, while the reported matrix is fixed-$N$
because sequential stopping biases effect-size estimates \cite{botlab}.

\subsection{Diagnostics}
\label{sec:metrics}

Alongside the payoff matrix, the lab records per-style diagnostics explaining \emph{why} a
style wins or loses: ask hit rate, turn retention, declare precision and yield, and a set of
\usff{}-specific metrics \cite{styles}: \code{concedeRate} (declares that gifted the set ---
the successor of the 48-card void rate, measuring a different event), \code{foreignDeclareRate}
(declares of sets the declarer held no card of), \code{declareLatency} (window-cycles between
a set becoming provable and the team declaring it, separating ``knew and waited'' from ``never
knew''), \code{raceLosses} (a teammate declared wrongly a set this seat could have declared
correctly), and \code{clinchDenials}. One shipped metric is explicitly excluded from style
comparisons: \code{hoardIndex}, a time-averaged count of sets held, does \emph{not} measure
hoarding --- a hoarding style lengthens games, adding late small-hand samples that drag the
mean down, a confound measured to be large enough to invert the sign \cite{styles}.

\subsection{The verdict rule}
\label{sec:verdict}

The decision rule was fixed in the experimental-design document before the tournament ran
\cite{botlab}: declare a superior style only if \emph{all four} hold ---

\begin{enumerate}
\item highest mean score; and
\item maximin $> 0.5$ (no losing matchup); and
\item cyclic energy of the matrix below $0.15$; and
\item exploitability not materially worse than its rivals'.
\end{enumerate}

If (2) or (3) fails, the published answer is ``no dominant style; here is the counter-graph
and the equilibrium mixture'' --- rendered as a first-class result. Cyclic energy is the
cyclic fraction of the Hodge decomposition of the duplicate-averaged antisymmetric payoff
matrix \cite{jiang2011,balduzzi2018}; the equilibrium mixture is the maximum-entropy Nash of
that matrix \cite{balduzzi2018}; $\alpha$-Rank supplies the evolutionary ranking
\cite{omidshafiei2019}. With 36 simultaneous cells, roughly $1.8$ would appear significant at
$\alpha = 0.05$ by chance alone, so every significance label --- including the individual
3-cycles behind criterion~3 --- is Benjamini--Hochberg corrected \cite{benjamini1995}.

\subsection{The decision-divergence instrument}
\label{sec:instrument}

Distinctness is not assumed from parameters; it is measured at the decision level. Two policy
vectors are compared by replaying real \usff{} games on the lab's own seed set and, at
\emph{every} decision point, handing both vectors the identical \code{SeatView} and seed and
comparing their two chosen \code{GameAction}s \cite{styles}. The instrument returns the exact
fraction of decisions on which the vectors differ; zero means the vectors are the same player.
It is used three ways below: to verify each style differs from the control
(Section~\ref{sec:distinct}), to expose the inert threshold axis (Section~\ref{sec:inert}),
and to detect that the first Hoarder implementation was byte-identical to Balanced
(Section~\ref{sec:hoarder}).

\subsection{Exploitability search}

For each style $i$, a best-response search over \code{StyleParams} deviations from $i$'s own
vector --- candidates drawn from a fixed knob ladder, each SPRT-gated, with an accepted
candidate re-measured by a fixed-$N$ 400-pair evaluation on seeds no search decision saw ---
reports $E(i) = \mathrm{SR}(\text{best response}, i) - 0.5$ \cite{botlab,styles}. $E$
distinguishes ``champion of this roster'' from ``actually strong'': a style can top the table
yet be maximally exploitable by a style nobody has written.

\subsection{Health gates and provenance}
\label{sec:provenance}

A run is void unless its health counters are clean. For the reported tournament
\cite{artifact}: \code{illegalActions} 0, \code{cappedGames} 0, \code{invariantViolations} 0,
\code{ties} 0, \code{voids} 0, \code{nonClinch} 0, \code{distinctSeeds} 4{,}300/4{,}300. The
artifact is committed at \code{src/lab/data/style-results.v2.json} with records digest
\code{0c3cb524a3534d4a}, engine commit \code{1667a1d} (recorded \code{1667a1d-dirty}, the
suffix reflecting generation from a modified working tree at that commit), and the SHA-256 of the
rules document (\code{b4c7fef7\ldots}) so that results become detectably stale if the rules
move. The full run --- 309{,}600 games plus analysis --- took roughly 16 minutes of wall
clock. Two matrix artifacts exist: \emph{v1} (engine commit \code{819eebb}, likewise recorded
\code{819eebb-dirty}), and \emph{v2},
re-run on the same seed set after the contained-book turn-pass was implemented
(Section~\ref{sec:containedpass}); all headline numbers in this paper are v2, with v1 cited
where the paired comparison is itself the result.

\section{Results}
\label{sec:results}

\subsection{The verdict: dominant --- Punter}

All four criteria of Section~\ref{sec:verdict} hold, and the artifact's computed verdict is
\emph{dominant} (Table~\ref{tab:criteria}). This is the criteria mechanism working as
designed, not a narrative choice: the same pipeline renders ``cyclic'' or ``inconclusive''
when the numbers say so.

\begin{table}[t]
\centering
\caption{The four dominance criteria, as computed in the committed artifact \cite{artifact}.}
\label{tab:criteria}
\small
\begin{tabular}{@{}clp{8.2cm}c@{}}
\toprule
\# & Criterion & Measured (matrix v2) & Pass \\
\midrule
1 & highest mean score & Punter $0.5829$, ahead of Blitz $0.5602$ (margin $0.0227$) & yes \\
2 & maximin $> 0.5$ & $0.5190$ (worst matchup: Balanced); CI$_{95}$ lower bound $0.5119$;
    worst cell significant after BH & yes \\
3 & cyclic energy $< 0.15$ & $0.0112$; zero significant 3-cycles after BH & yes \\
4 & exploitability not materially worse & $E(\text{Punter}) = 0.0025$ vs.\ rivals' median
    $0.0444$ & yes \\
\bottomrule
\end{tabular}
\end{table}

\subsection{The ranking, five ways}

Table~\ref{tab:ranking} shows the five ranking systems side by side. They agree almost
everywhere --- which is itself evidence of transitivity. The mean-score, Bradley--Terry and
Hodge-rating orderings are identical position for position; $\alpha$-Rank differs only in
swapping two adjacent pairs deep in the table: Hoarder/Archivist (mean-score gap $0.0083$,
about $3.6$ combined standard errors of the means, though the pair's head-to-head cell is
itself a statistical tie) and Scout/Turtle (gap $0.0038$, about $1.6$ combined standard errors
--- a statistical tie in mean score). Nash averaging places all its mass
on Punter, and $\alpha$-Rank concentrates on it likewise --- the degenerate mixtures that a
transitive matrix with a strict maximin winner forces.

\begin{table}[t]
\centering
\caption{Ranking systems side by side (matrix v2, 4{,}300 pairs per cell)
\cite{artifact}. B--T is the Bradley--Terry rating in Elo-like units.}
\label{tab:ranking}
\footnotesize
\begin{tabular}{@{}lcccccrc@{}}
\toprule
Style & Mean score & 95\% CI & Maximin & Worst vs. & Nash mass & B--T & $\alpha$-Rank \\
\midrule
Punter    & 0.5829 & [0.5798, 0.5860] & 0.5190 & Balanced & 1.000 & $52.2$ & 1 \\
Blitz     & 0.5602 & [0.5570, 0.5633] & 0.4802 & Punter   & 0.000 & $37.8$ & 2 \\
Balanced  & 0.5581 & [0.5551, 0.5612] & 0.4810 & Punter   & 0.000 & $36.5$ & 3 \\
Banker    & 0.5285 & [0.5252, 0.5318] & 0.4531 & Punter   & 0.000 & $17.9$ & 4 \\
Ghost     & 0.4928 & [0.4894, 0.4962] & 0.4133 & Punter   & 0.000 & $-4.4$ & 5 \\
Hoarder   & 0.4885 & [0.4854, 0.4917] & 0.4205 & Punter   & 0.000 & $-7.1$ & 7 \\
Archivist & 0.4802 & [0.4769, 0.4835] & 0.4074 & Punter   & 0.000 & $-12.3$ & 6 \\
Scout     & 0.4063 & [0.4030, 0.4096] & 0.3397 & Punter   & 0.000 & $-59.1$ & 9 \\
Turtle    & 0.4025 & [0.3992, 0.4058] & 0.3419 & Punter   & 0.000 & $-61.5$ & 8 \\
\bottomrule
\end{tabular}
\end{table}

The worst-vs.\ column carries a structural fact: every style's worst matchup is Punter. Punter
itself has no losing matchup; its narrowest, against Balanced and Blitz, are two cells that
are statistically tied with each other ($0.5190$ and $0.5198$) --- indeed, between matrix v1
and v2 the identity of Punter's worst opponent flipped between these two, a re-ordering of
tied cells rather than a new counter \cite{styles}. Punter's full row is
Table~\ref{tab:punterrow}: every one of its eight cells is above $0.5$ and significant after
BH correction.

\begin{table}[t]
\centering
\caption{Punter's matchup row (score from Punter's side; all cells BH-significant)
\cite{artifact}.}
\label{tab:punterrow}
\footnotesize
\begin{tabular}{@{}lcc@{}}
\toprule
Opponent & Punter's score & SE \\
\midrule
Balanced  & 0.5190 & 0.0036 \\
Blitz     & 0.5198 & 0.0042 \\
Banker    & 0.5469 & 0.0048 \\
Hoarder   & 0.5795 & 0.0042 \\
Ghost     & 0.5867 & 0.0049 \\
Archivist & 0.5926 & 0.0047 \\
Turtle    & 0.6581 & 0.0046 \\
Scout     & 0.6603 & 0.0046 \\
\bottomrule
\end{tabular}
\end{table}

\subsection{Matrix structure: almost no cycles}

Thirty-three of the 36 cells are significant after BH correction. The three that are not are
exactly the near-ties one would hope a well-powered design exposes as near-ties:
Balanced--Blitz ($0.4962$, $q = 0.33$), Hoarder--Archivist ($0.5074$, $q = 0.12$) and
Ghost--Archivist ($0.4944$, $q = 0.26$). The Hodge decomposition puts the cyclic energy at
$0.0112$ --- about $1.1\%$ of a matrix whose verdict threshold is $15\%$. Exactly one directed
3-cycle exists in point estimates (Hoarder $>$ Archivist $>$ Ghost $>$ Hoarder, minimum edge
$0.5056$), and it is not significant after BH ($q = 0.27$) \cite{styles,artifact}. The
counter-structure that the laboratory was built to detect, had it existed, is measurably
absent at this precision: the styles form a ladder, not a wheel.

\subsection{Why Punter wins}
\label{sec:why}

Table~\ref{tab:diag} pools each style's diagnostics over its eight cells. Punter's profile is
volume at near-control precision: it makes the most declares per game ($4.08$ against the
control's $3.99$) and converts the most sets ($4.11$ against $4.02$), while its declare
precision gives up only $0.2$ points against the control ($0.9884$ vs.\ $0.9901$) and its
concession rate rises only from $0.0099$ to $0.0116$. It pays measurable costs --- more chosen
concessions ($0.0066$ of its declares against the control's $0.0042$) and more race losses
($0.045$ per game against $0.037$) --- and the two-point swing of
Section~\ref{sec:signflip} prices those costs; they are simply smaller than the value of the
extra conversions.

The mechanism behind that profile is \emph{not} the one the style's label advertises.
Section~\ref{sec:inert} shows that Punter's low declare threshold ($0.775$) selects exactly
the same declares as the control's $0.90$: the axis is inert. What differs is the ask policy
--- \code{gambleBonus} 25 steers Punter toward asks that would complete a set, and
\code{declareMaxUncertain} 2 lets it tolerate a second guessed card where every other style
tolerates at most one. Punter wins as the style that hunts completions and accepts the
priced risk of the gift rule, not as the style that declares on thinner evidence.

At the other end, the table explains the floor. Turtle's own-hand-only gate leaves it holding
sets it cannot bank; its declares are forced ones over a third of the time ($0.367$ against
the control's $0.069$), its precision collapses to $0.9456$, its concession rate is five times
the control's, and it converts barely $2.94$ sets per game. Scout buys information with its
ask weights ($\code{wHit}$ 55, $\code{wNarrow}$ 40) and pays for it in material: the
second-lowest hit rate and the fewest conversions of the non-Turtle styles. Ghost is the
mirror image --- the roster's \emph{highest} hit rate ($0.6025$), bought by leak-protection
that steers its asks away from informative sets --- and still finishes below $0.5$: under
\usff{}, protecting the read is worth less than the tempo it costs. The Hoarder's line is
read in Section~\ref{sec:hoarder}.

\begin{table}[t]
\centering
\caption{Per-style diagnostics, pooled over each style's eight cells (matrix v2)
\cite{artifact}. ``Hit'' = ask hit rate; ``Prec'' = declare precision; ``Conc'' =
concession rate (share of the style's declares gifted to the opponents); ``Frgn'' =
foreign-declare rate; ``Lat'' = declare latency (window-cycles); ``Race'' = race losses per
game; ``Sets'' and ``Decl'' = sets converted and declares made per game; ``Forced'' = share
of declares that were compelled.}
\label{tab:diag}
\footnotesize
\begin{tabular}{@{}lccccccccc@{}}
\toprule
Style & Hit & Prec & Conc & Frgn & Lat & Race & Sets & Decl & Forced \\
\midrule
Punter    & 0.564 & 0.988 & 0.012 & 0.013 & 22.8 & 0.045 & 4.11 & 4.08 & 0.066 \\
Blitz     & 0.564 & 0.990 & 0.010 & 0.014 & 22.9 & 0.036 & 4.03 & 4.00 & 0.067 \\
Balanced  & 0.568 & 0.990 & 0.010 & 0.014 & 23.2 & 0.037 & 4.02 & 3.99 & 0.069 \\
Banker    & 0.563 & 0.990 & 0.010 & 0.014 & 22.8 & 0.035 & 3.93 & 3.90 & 0.072 \\
Ghost     & 0.603 & 0.989 & 0.011 & 0.010 & 24.3 & 0.035 & 3.77 & 3.74 & 0.099 \\
Hoarder   & 0.545 & 0.975 & 0.025 & 0.253 & 30.6 & 0.088 & 3.66 & 3.69 & 0.148 \\
Archivist & 0.556 & 0.967 & 0.033 & 0.007 & 19.8 & 0.091 & 3.76 & 3.79 & 0.082 \\
Scout     & 0.532 & 0.957 & 0.043 & 0.007 & 17.2 & 0.088 & 3.49 & 3.52 & 0.080 \\
Turtle    & 0.523 & 0.946 & 0.054 & 0.060 & 63.4 & 0.162 & 2.94 & 3.00 & 0.367 \\
\bottomrule
\end{tabular}
\end{table}

\subsection{Exploitability}
\label{sec:exploit}

Table~\ref{tab:exploit} reports $E(i)$ for each style with the accepted best-response
deviation. Punter's $E$ of $0.0025$ --- its best measured counter being merely to randomize
the miss target --- confirms it is not a paper champion of this particular roster. The
ordering tracks the ranking: the exploitable styles are the ones already at the bottom, with
Turtle's single-knob repair (turn off the own-hand-only gate) worth a full $0.15$.

Two honest cautions accompany the table \cite{styles}. $E$ is a maximum over a stochastic
search whose final evaluations carry standard errors of $0.008$--$0.017$ and which only
accepts moves worth at least ${\approx}0.03$, so small entries and small differences between
runs are not meaningful. And the search's knob ladder predates the contained-book turn-pass,
so that one axis was probed separately by the same SPRT protocol: no candidate was accepted
for any style at any appetite --- consistent with Section~\ref{sec:containedpass}'s finding
that the mechanism is worth nothing.

\begin{table}[t]
\centering
\caption{Exploitability, matrix v2: $E(i)$ and the accepted best-response deviation from the
style's own vector \cite{artifact,styles}.}
\label{tab:exploit}
\footnotesize
\begin{tabular}{@{}lcl@{}}
\toprule
Style & $E$ & Accepted deviation \\
\midrule
Balanced  & 0.0000 & none (mirror) \\
Banker    & 0.0012 & \code{gambleBonus} 15 \\
Punter    & 0.0025 & \code{missTarget} \code{random} \\
Blitz     & 0.0175 & \code{leakEpsilon} 3 \\
Archivist & 0.0325 & \code{wHit} 100 \\
Hoarder   & 0.0563 & \code{hoardBooks} 0, \code{minHandSize} 0 --- \emph{stop hoarding} \\
Ghost     & 0.0775 & \code{wProgress} 30, \code{gambleBonus} 30 \\
Scout     & 0.1112 & \code{wHit} 90, \code{wNarrow} 0, \code{minHitP} 0.45 \\
Turtle    & 0.1500 & \code{declareOnlyOwnHand} off \\
\bottomrule
\end{tabular}
\end{table}

\section{Caveat I: the declare-threshold axis is inert}
\label{sec:inert}

This section states the finding, its mechanism, and its consequences for reading the results;
the full treatment --- instrument, attribution, repairs, and the methodological argument --- is
a companion paper published alongside this one \cite{inertpaper}.

\subsection{The styles are distinct\ldots}
\label{sec:distinct}

Run through the decision-divergence instrument (Section~\ref{sec:instrument}) against the
Balanced control, every non-control style differs on a measurable fraction of decisions
\cite{styles}:

\begin{center}
\footnotesize
\begin{tabular}{@{}lcccccccc@{}}
\toprule
 & Scout & Ghost & Archivist & Banker & Turtle & Punter & Hoarder & Blitz \\
\midrule
\% of decisions differing & 2.89 & 2.15 & 1.66 & 1.19 & 0.83 & 0.47 & 0.47 & 0.39 \\
\bottomrule
\end{tabular}
\end{center}

\noindent
Between $0.39\%$ and $2.89\%$ of decisions may look small, but Literature is a game where a
single redirected ask changes who holds the turn and what the whole table learns; the
tournament shows these divergence rates compounding into mean-score gaps of up to 18 points.

\subsection{\ldots but not along the axis their labels advertise}

\code{declareThreshold} is the knob on which the roster's aggressive-to-conservative story is
built. Varying it \emph{alone} on the Balanced control, over 40 \usff{} games with every
decision compared under the instrument, gives Table~\ref{tab:sweep}: the divergence collapses
to exactly zero at $0.775$ and stays zero at every measured point through $0.975$ (the pilot
sweep's zero extends through threshold $1$) --- and \emph{every roster style's shipped
threshold sits at $0.775$ or above}, the shipped span being $[0.775, 0.98]$. The axis is dead
for all nine of them \cite{styles}.

\begin{table}[t]
\centering
\caption{Decision divergence when only \code{declareThreshold} varies on the Balanced
control (40 games) \cite{styles}. All roster styles ship thresholds in $[0.775, 0.98]$.}
\label{tab:sweep}
\footnotesize
\begin{tabular}{@{}lccccccccc@{}}
\toprule
\code{declareThreshold} & 0.50 & 0.65 & 0.70 & 0.75 & \textbf{0.775} & 0.85 & 0.90 & 0.95 & 0.975 \\
\midrule
divergent decisions & 350 & 48 & 29 & 19 & \textbf{0} & 0 & 0 & 0 & 0 \\
\% of decisions & 1.297 & 0.178 & 0.108 & 0.070 & \textbf{0.000} & 0.000 & 0.000 & 0.000 & 0.000 \\
\bottomrule
\end{tabular}
\end{table}

The mechanism is not an unreachable code path --- Balanced makes 171 speculative declares to
138 certain ones over those same 40 games, so the thresholded branch fires constantly. (An
earlier audit, described in Section~\ref{sec:hoarderaudit}, had read a related zero as the
branch never being reached; the direct sweep corrected that reading.) The measured mechanism
is that the shared inference engine's confidence estimates are \textbf{bimodal}: a candidate
declaration is either well above ${\sim}0.975$ or well below ${\sim}0.775$, with almost
nothing in between, so any threshold inside that band selects the identical set of declares
\cite{styles}.

Three consequences follow, and each is stated wherever the ranking is published. First,
Punter did not win by declaring at $0.775$ rather than $0.90$; it won through
\code{gambleBonus} and \code{declareMaxUncertain} (Section~\ref{sec:why}). Second, any
description of this roster as a spectrum from aggressive to conservative \emph{declaring}
describes a knob that does not fire; the labels survive only as names for appetites. Third,
the $(1+t)/2$ appetite mapping of Section~\ref{sec:signflip}, while arithmetically right, is
currently unfalsifiable on this engine: every value it can produce in the plausible range
yields identical play \cite{styles}.

\section{Caveat II: the Hoarder --- costs measured, benefit implemented, still worthless}
\label{sec:hoarder}

The Hoarder is the project's originally named style, and it is mechanically grounded: holding
at least one card of a set is the \emph{only} licence to ask into it, so declaring a set
revokes your own asking rights there. Its story across v0.5 is the paper's second negative
result, and it runs in three acts: an implementation that turned out to be the control in
disguise; an audited implementation that expresses the thesis and loses by it; and the
style's theorized benefit, implemented and measured to be worth nothing. The third act's
mechanism --- the contained-book turn-pass, its theorem, trigger, and measurements --- is
treated in full in a second companion paper published alongside this one \cite{containedpaper}.

\subsection{Act one: the inert implementation}
\label{sec:hoarderaudit}

In the 36-cell pilot (200 pairs per cell), the Hoarder's mean score was $0.5484$ --- identical
to Balanced's in every digit --- its rows of the payoff and diagnostics tables were
byte-identical to the control's, and the head-to-head cell measured $0.5000$ with a standard
error of exactly zero. Rather than infer, the identity was instrumented: over 400 games and
275{,}380 decision points, the Hoarder and Balanced vectors, handed identical views and seeds,
chose \emph{zero} different actions \cite{styles}.

Attributing the zero one parameter at a time (each of the Hoarder's seven behavioural deltas
applied to Balanced alone, over 76{,}972 decisions) returned zero for all seven. The audit's
first finding was structural: the two hoard knobs gated only the speculative-declare branch,
on the argument that refusing \emph{free} points to protect an ask-licence is a giveaway
rather than optionality --- but within the measured population that branch never selected
differently across the whole threshold band (the flatness later explained by
Section~\ref{sec:inert}'s bimodality), so the gate guarded a distinction without a
difference, and with it the entire style. The fix moved the hoard gate in front of
\emph{every} declare the style is free to refuse --- speculative, certain, or wholly in hand
--- while never touching the compulsory \code{MUST\_DECLARE} positions. The two knob values
were then derived from the rule set rather than tuned: \code{minHandSize} 2 is the smallest
hand that survives one involuntary strip (hits are compulsory) and still holds a licence,
the only hand-size discontinuity in the rules being at zero cards; \code{hoardBooks} 3 is the
style's stated appetite, corrected to count only licences with a legal ask behind them
(a set the seat would hold all six of grants a licence with nothing legal to name)
\cite{styles}.

\subsection{Act two: expressible, and wrong}

After the audit the same instrument measures $9{,}315$ divergent decisions in $289{,}560$
($3.22\%$): the style exists. Every diagnostic moves in the direction the thesis predicts
(Table~\ref{tab:hoarderpilot}): it stays in the asking game longer (dropout rate $0.521 \to
0.443$), waits $35\%$ longer to bank (declare latency $22.90 \to 31.02$ window-cycles), and --- because the sets
it does declare are disproportionately the ones that cost it no licences --- its
foreign-declare rate rises thirteen-fold, the project owner's ``memorize half-suits you do not
own'' strategy arriving as a side effect of hoarding. And it pays for all of it: the
head-to-head against Balanced moves from $0.5000 \pm 0.0000$ to $0.5775 \pm 0.0159$
\emph{against} the Hoarder, and its pilot mean falls from $0.5484$ to $0.4738$. Ablated on
identical deals, \code{hoardBooks} 3 alone reproduces the entire effect ($0.4225$ to the
Hoarder, identical to both knobs together, because any declare leaving under two cards also
leaves under three licences); with both knobs off the vector returns to $0.5000 \pm 0.0000$
--- the inertness finding restated as an outcome \cite{styles}.

\begin{table}[t]
\centering
\caption{The Hoarder before and after the audit --- pilot diagnostics pooled over its eight
cells \cite{styles}. The inert version is byte-identical to the control. Two Balanced figures
coexist in the source: $0.5484$ is the pooled value in the identity comparison that flagged the
inert Hoarder, while this table's Balanced column pools Balanced's own eight cells ($0.5581$).
The load-bearing identity is byte-identical play --- zero divergent decisions in 275{,}380 ---
and the table reports both figures rather than resolving what the source does not.}
\label{tab:hoarderpilot}
\footnotesize
\begin{tabular}{@{}lccc@{}}
\toprule
Diagnostic (pilot) & Balanced & Hoarder, inert & Hoarder, audited \\
\midrule
mean score            & 0.5581 & 0.5484 (= Balanced) & 0.4738 \\
foreign-declare rate  & 0.018  & 0.018  & 0.247 \\
declare latency       & 23.39  & 22.90  & 31.02 \\
race losses / game    & 0.044  & 0.046  & 0.091 \\
declares / game       & 4.01   & 3.98   & 3.67 \\
forced-declare share  & 0.074  & 0.076  & 0.153 \\
concession rate       & 0.011  & 0.012  & 0.025 \\
dropout rate          & 0.539  & 0.521  & 0.443 \\
ask hit rate          & 0.568  & 0.569  & 0.545 \\
\bottomrule
\end{tabular}
\end{table}

In the full-precision matrices the audited Hoarder finishes 6th of 9 (the pilot's $0.4738$
figure corresponds to a 7th-place pilot finish; both full-precision runs place it 6th
\cite{styles}). At this point an objection remained open, and the project's containment
analysis \cite{containment} had sharpened it: the verdict was measured on an implementation
that paid hoarding's full costs while collecting none of its benefit.

\subsection{Act three: the benefit, implemented --- the contained-book turn-pass}
\label{sec:containedpass}

Call a set \emph{contained} when all six of its cards sit with one team. Containment is an
absorbing state with measured consequences \cite{containment}: no opponent can legally ask
into the set (asking requires holding a card of it); an opponent declaring it necessarily
assigns to their own team and thus gifts it straight back; a holder may still ask into it ---
a legal, \emph{guaranteed} miss that passes the turn to a chosen opponent; the miss moves no
cards, so the licence is never consumed; and declaring the set destroys the move. Claiming a
contained set therefore has no defensive value, and leaving it unclaimed preserves a
repeatable, aimable turn-pass --- advice independently attested by Develin's strategy text
\cite{develin}. This is precisely the benefit a hoarder's unclaimed sets should be earning.

v0.5 implements the turn-pass as a policy option with a derived, unturned trigger. Writing
$p^\ast$ for the hit probability of the ask the style would otherwise play, the pass is taken
iff
\[
p^\ast \;<\; \frac{a \cdot \mathrm{gain} - \mathrm{infoCost}}{\mathrm{tempo}},
\]
with every term in units of cards and read off the public log: the expected turn length
$E = \mathrm{hits}/\max(1, \mathrm{misses})$ prices a conceded turn at $E \cdot n_t/\bar n$
for an opponent holding $n_t$ cards (mean live hand $\bar n$), $\mathrm{gain}$ is the spread
between the opponent the ordinary ask would concede to and the opponent the style's
\code{missTarget} prefers, $\mathrm{tempo} = 1 + E \cdot n_o/\bar n$ is the swing between
hitting and missing, and $\mathrm{infoCost}$ is zero on a reuse of an already-published card
and $1/U$ otherwise ($U$ the count of live cards not yet publicly placed) \cite{styles}. The
appetite $a$ is the expected number of uses before the set is banked: 1 for eight of the nine
styles, and $1.33$ for the Hoarder --- the one style whose delay had actually been measured,
the number read off its own declare-latency ratio. A certain hit is never displaced, and the
policy reuses one card, so $94.8$--$97.6\%$ of uses by the heavy users cost no information at
all. One honest check on the derivation: measured uses per contained book order the styles
exactly as the appetite argument predicts, but the Hoarder-to-Balanced ratio comes out $1.18$,
not the $1.33$ the latency ratio implied --- right in sign, roughly right in size, and
deliberately \emph{not} re-fitted, since tuning an appetite to the behaviour it produces would
be fitting to the outcome \cite{styles}.

\begin{table}[t]
\centering
\caption{The contained-book turn-pass: opportunity, firing, and value \cite{styles}.
Opportunity and firing rates are from 100 mirror games per style; the score is the style
\emph{with} the mechanism against the identical style without it, over 1{,}200 duplicate
pairs (2{,}400 games); ``changed'' counts games the mechanism actually touched.}
\label{tab:contained}
\footnotesize
\begin{tabular}{@{}lcccccc@{}}
\toprule
Style & Appetite & Opportunity \% & Fired / opp.\ \% & Divergence \% & Score (with vs.\ without) & Changed \\
\midrule
Balanced  & 1    & 3.89  & 7.7  & 0.041 & $0.5046 \pm 0.0034$ & 239 \\
Blitz     & 1    & 5.05  & 0.0  & 0.000 & $0.5000 \pm 0.0000$ & 0 \\
Punter    & 1    & 3.65  & 3.8  & 0.019 & $0.4971 \pm 0.0029$ & 142 \\
Banker    & 1    & 5.04  & 13.9 & 0.098 & $0.5021 \pm 0.0035$ & 281 \\
Turtle    & 1    & 53.86 & 7.3  & 0.558 & $0.5008 \pm 0.0065$ & 1{,}268 \\
Hoarder   & 1.33 & 14.95 & 12.4 & 0.258 & $0.4938 \pm 0.0054$ & 915 \\
Scout     & 1    & 0.38  & 0.0  & 0.000 & $0.5000 \pm 0.0000$ & 0 \\
Ghost     & 1    & 1.31  & 16.0 & 0.029 & $0.4958 \pm 0.0027$ & 156 \\
Archivist & 1    & 1.01  & 4.8  & 0.007 & $0.4996 \pm 0.0004$ & 11 \\
\bottomrule
\end{tabular}
\end{table}

Table~\ref{tab:contained} reports what the mechanism does. Three readings. First, the
opportunity is a property of the \emph{declare} policy: a contained set exists at $54\%$ of
the Turtle's ask decisions (it never banks sets split across teammates) and $0.4\%$ of the
Scout's (it pins holders early, and a pinned contained set is banked immediately) --- and the
Hoarder sits second at $14.95\%$, nearly four times the control's $3.89\%$, which is exactly
the benefit the objection said it was paying for and not collecting. Second, even where the
licence exists it is mostly refused ($3.8$--$16\%$ taken), always on the derivation's
$\mathrm{gain} \le 0$ arm: the ordinary ask was already conceding the turn where the style
wanted it. Blitz needs no exception --- its \code{missTarget} \code{most} makes the aiming
gain non-positive, so the mechanism turns itself off. Third, and decisively: \textbf{no style
is two standard errors from $0.5$}. The largest deviation is $1.6$ SE (Ghost), the two heavy
users land on and slightly below the null (Turtle $0.5008 \pm 0.0065$, Hoarder $0.4938 \pm
0.0054$), and the styles that never fire return exactly $0.5000 \pm 0.0000$ over 2{,}400
games each --- the with- and without-vectors are the same player. \emph{The measured value of
the contained-book turn-pass, as a policy option under \usff{} against this roster, is zero.}
The arithmetic says why: with measured $E \approx 1$, the whole spread between the best and
worst opponent to concede to is on the order of one card, against which the displaced ask was
worth $p^\ast$ of a card plus a retained turn. The move is real, and it is small.

\subsection{The verdict on the strategy}

The full 36-cell matrix was then re-run with the mechanism live, at the same precision on the
same seed set, so v1 and v2 differ by the policy change and nothing else. Nothing that
matters moved (Table~\ref{tab:v1v2}): the verdict, the complete ranking, the Nash and
$\alpha$-Rank concentrations, and every rank are identical; no cell moved significantly (the
largest paired per-deal difference, Blitz--Ghost at $+0.0030 \pm 0.0015$, $t = 2.07$, does
not survive BH over 36 comparisons --- one $|t| > 2$ in 36 is what chance produces); cyclic
energy went $0.0097 \to 0.0112$ with the same single non-significant 3-cycle; and the
significant-cell count's $34 \to 33$ change is entirely the Hoarder--Archivist cell drifting
out of significance ($q\ 0.036 \to 0.118$) \cite{styles}.

\begin{table}[t]
\centering
\caption{Matrix v1 vs.\ v2: per-style mean score, paired over the same 4{,}300 deals
\cite{styles}. Every delta is within $1.5$ SE of zero; every rank is unchanged.}
\label{tab:v1v2}
\footnotesize
\begin{tabular}{@{}lcccccc@{}}
\toprule
Style & v1 & v2 & $\Delta$ & SE($\Delta$) & $t$ & Rank \\
\midrule
Punter    & 0.5829 & 0.5829 & $-0.0001$ & 0.0009 & $-0.07$ & $1 \to 1$ \\
Blitz     & 0.5594 & 0.5602 & $+0.0008$ & 0.0008 & $+0.94$ & $2 \to 2$ \\
Balanced  & 0.5584 & 0.5581 & $-0.0003$ & 0.0009 & $-0.34$ & $3 \to 3$ \\
Banker    & 0.5274 & 0.5285 & $+0.0011$ & 0.0011 & $+1.00$ & $4 \to 4$ \\
Ghost     & 0.4928 & 0.4928 & $+0.0000$ & 0.0010 & $+0.03$ & $5 \to 5$ \\
Hoarder   & 0.4904 & 0.4885 & $-0.0018$ & 0.0014 & $-1.35$ & $6 \to 6$ \\
Archivist & 0.4791 & 0.4802 & $+0.0011$ & 0.0008 & $+1.47$ & $7 \to 7$ \\
Scout     & 0.4064 & 0.4063 & $-0.0001$ & 0.0007 & $-0.12$ & $8 \to 8$ \\
Turtle    & 0.4032 & 0.4025 & $-0.0007$ & 0.0019 & $-0.38$ & $9 \to 9$ \\
\bottomrule
\end{tabular}
\end{table}

For the Hoarder specifically, the prediction that the missing benefit explained its finish
\emph{did not hold}. Its paired delta is $-0.0018 \pm 0.0014$ ($t = -1.35$) --- the point
estimate went \emph{down}, indistinguishably from zero; not one of its eight cells moved
significantly, and its one marginal cell moved against significance. The sharpest form of the
result comes from the exploitability search, free to pick any single knob: the best response
to the Hoarder-with-the-benefit is \code{hoardBooks} 0, \code{minHandSize} 0 --- \emph{stop
hoarding} --- worth $0.5563 \pm 0.0124$ (Table~\ref{tab:exploit}). One number from the
follow-up single-axis search deserves recording precisely because half of it would have
looked like a finding: probing \code{containedPass} $= 0$ against the Hoarder, the SPRT
search seeds said the candidate was \emph{worse} by $0.020 \pm 0.018$ while the held-out
evaluation said \emph{better} by $0.0275 \pm 0.0098$; two contradictory effects of that size
mean neither is real, and the honest summary is noise \cite{styles}.

So the conclusion is about the strategy, not the implementation. The benefit the containment
analysis identified is real, reachable, and fires for the derived reasons --- and it is worth
nothing at this roster's level of play. The verdict stands with no unimplemented benefit
behind it as an excuse: \textbf{optionality under \usff{} is real, purchasable, and priced
above its worth}, and the two rules doing the pricing are the ones
Section~\ref{sec:signflip} flagged --- the gift rule makes a teammate's wrong declare of your
hoarded set a two-point event (the Hoarder's race losses double), and the right of cardless
players to keep declaring means the elimination the hoard defends against is only elimination
from the asking half of the game. What this does \emph{not} settle: the trigger competes
against the roster's tuned ask policies, the systematic re-tuning protocol has never been run
(Section~\ref{sec:threats}), and the second containment mechanism --- holding contained sets
by default in the \emph{declare} policy, which the opportunity-rate column suggests is where
the leverage actually lives --- is untouched in v0.5 \cite{styles}.

\section{Threats to validity}
\label{sec:threats}

Stated as designed-for or open, per the laboratory's own checklist \cite{botlab}.

\begin{enumerate}
\item \textbf{Style as handicap.} Designed against: one shared full-strength inference engine;
  distinctness verified at the decision level. \emph{Open:} the specified skill ablation ---
  re-running every style over medium-strength inference to detect style-times-skill
  interactions --- has not been run for this roster.
\item \textbf{Self-play convention lock-in.} A style that wins this tournament may be a style
  its own roster has learned to read \cite{hu2020}. \emph{Open:} the holdout roster and the
  cross-play protocol against an independently written bot are specified but not yet run; the
  artifact's cross-play block is empty.
\item \textbf{Intransitivity misread as noise.} Designed against: explicit cyclic-energy
  metric, BH-corrected cycle enumeration, and a verdict rule with a first-class cyclic
  outcome.
\item \textbf{Tie-blindness.} Moot under \usff{}: ties are impossible
  (Theorem~\ref{thm:tie}) and the artifact asserts zero across 309{,}600 games.
\item \textbf{Passive-style stalling.} The step cap is instrumented, the stall-breaker is
  global (a per-style stall rule would be a hidden style parameter), and
  \code{cappedGames} $= 0$ in the reported run.
\item \textbf{Determinism mistaken for sampling.} Re-running a seed reproduces the game
  byte-for-byte; the harness asserts 4{,}300 distinct seeds.
\item \textbf{Multiplicity.} All 36 cells and all cycle tests are BH-corrected.
\item \textbf{Conditional on the rule set.} Every conclusion is conditional on \usff{} as
  pinned by the artifact's rules hash. The sign-flip analysis is itself evidence that results
  do not transfer across dialects; nothing here is claimed for the 48-card game.
\item \textbf{An untuned roster.} The styles are evaluated \emph{as specified}: the
  SPRT-based re-tuning protocol (each style tuned against the control) has never been run, so
  every per-style result --- including both negative results --- is about the specified
  vector, not the best version of the style that exists. A roster with weaker ask policies
  would, in particular, leave more room for the turn-pass of
  Section~\ref{sec:containedpass}.
\item \textbf{Uniform teams only.} All cells are pure teams of three. The mixed-composition
  tier --- including the hypothesis that the Archivist is weak in mirrors but strong as a
  single seat --- is specified and unexplored; the artifact's teams block is empty.
\end{enumerate}

\section{Conclusion}
\label{sec:conclusion}

Under the \usff{} dialect, with style isolated from skill and the payoff matrix decomposed
rather than assumed transitive, the question ``is there a superior play style?'' has a
measured answer: yes --- at this roster, at this level of play, the completion-chasing Punter
dominates by all four pre-stated criteria, the matrix carries almost no cyclic energy, and
the anticipated counter-structure is absent. The two results that qualify the headline are as
load-bearing as the headline itself: the declare-threshold axis the styles are named after is
inert across the roster's entire range, because the shared engine's confidence estimates are
bimodal; and the hoarding thesis fails end-to-end, its theorized benefit surviving
implementation only to be measured at zero. Honest instrumentation did not merely decorate
these findings --- the decision-divergence instrument caught a control masquerading as a
style, and the paired v1/v2 re-run kept a nothing-result from being narrated into a
something-result.

The natural next question is dynamic rather than static: whether an agent that identifies its
opponents' styles online and adapts its own can beat the fixed Punter. That is the subject of
separate, forthcoming work --- distinct from the two companion papers on the inert axis and
the contained book published alongside this one \cite{inertpaper,containedpaper}.

\medskip
\noindent\textbf{Reproducibility.} The engine is deterministic; the committed artifact
\cite{artifact} carries the seed set, per-cell records digest, engine commit and rules hash;
and every number in this paper is either in that artifact or in the repository's measurement
documents \cite{rules,styles,botlab,containment}.

\begin{thebibliography}{99}
\small

\bibitem{balduzzi2018}
D.~Balduzzi, K.~Tuyls, J.~P\'erolat, and T.~Graepel.
\newblock Re-evaluating evaluation.
\newblock In \emph{Advances in Neural Information Processing Systems (NeurIPS)}, 2018.

\bibitem{omidshafiei2019}
S.~Omidshafiei, C.~Papadimitriou, G.~Piliouras, K.~Tuyls, M.~Rowland, J.-B.~Lespiau,
W.~M.~Czarnecki, M.~Lanctot, J.~P\'erolat, and R.~Munos.
\newblock $\alpha$-Rank: Multi-agent evaluation by evolution.
\newblock \emph{Scientific Reports}, 9:9937, 2019.

\bibitem{jiang2011}
X.~Jiang, L.-H.~Lim, Y.~Yao, and Y.~Ye.
\newblock Statistical ranking and combinatorial Hodge theory.
\newblock \emph{Mathematical Programming}, 127(1):203--244, 2011.

\bibitem{bradleyterry1952}
R.~A.~Bradley and M.~E.~Terry.
\newblock Rank analysis of incomplete block designs: I. The method of paired comparisons.
\newblock \emph{Biometrika}, 39(3/4):324--345, 1952.

\bibitem{wald1945}
A.~Wald.
\newblock Sequential tests of statistical hypotheses.
\newblock \emph{Annals of Mathematical Statistics}, 16(2):117--186, 1945.

\bibitem{benjamini1995}
Y.~Benjamini and Y.~Hochberg.
\newblock Controlling the false discovery rate: a practical and powerful approach to
multiple testing.
\newblock \emph{Journal of the Royal Statistical Society, Series B}, 57(1):289--300, 1995.

\bibitem{burch2018}
N.~Burch, M.~Schmid, M.~Morav\v{c}\'{\i}k, D.~Morrill, and M.~Bowling.
\newblock AIVAT: A new variance reduction technique for agent evaluation in imperfect
information games.
\newblock In \emph{Proceedings of the AAAI Conference on Artificial Intelligence}, 2018.

\bibitem{cowling2012}
P.~I.~Cowling, E.~J.~Powley, and D.~Whitehouse.
\newblock Information set Monte Carlo tree search.
\newblock \emph{IEEE Transactions on Computational Intelligence and AI in Games},
4(2):120--143, 2012.

\bibitem{long2010}
J.~Long, N.~R.~Sturtevant, M.~Buro, and T.~Furtak.
\newblock Understanding the success of perfect information Monte Carlo sampling in game
tree search.
\newblock In \emph{Proceedings of the AAAI Conference on Artificial Intelligence}, 2010.

\bibitem{celli2018}
A.~Celli and N.~Gatti.
\newblock Computational results for extensive-form adversarial team games.
\newblock In \emph{Proceedings of the AAAI Conference on Artificial Intelligence}, 2018.

\bibitem{hu2020}
H.~Hu, A.~Lerer, A.~Peysakhovich, and J.~Foerster.
\newblock ``Other-play'' for zero-shot coordination.
\newblock In \emph{Proceedings of the International Conference on Machine Learning (ICML)},
2020.

\bibitem{bard2020}
N.~Bard et al.
\newblock The Hanabi challenge: A new frontier for AI research.
\newblock \emph{Artificial Intelligence}, 280:103216, 2020.

\bibitem{neller2018}
T.~W.~Neller and Z.~Luo.
\newblock Mixed logical and probabilistic reasoning in the game of Clue.
\newblock \emph{ICGA Journal}, 40(4):406--416, 2018.

\bibitem{develin}
M.~Develin.
\newblock \emph{The Ten Best Card Games You've Never Heard Of}, chapter~9 (Canadian Fish).
\newblock Self-published manuscript, \url{http://www.bantha.org/~develin/cardgames.html}.

\bibitem{pagat}
J.~McLeod (ed.).
\newblock Literature.
\newblock Card game rules at pagat.com,
\url{https://www.pagat.com/quartet/literature.html}.

\bibitem{somani}
N.~Somani.
\newblock literature: a Q-learning Literature bot (four-player; no reported win rates).
\newblock GitHub repository, \url{https://github.com/neelsomani/literature}, 2018.

\bibitem{playstyles}
FishAI project.
\newblock PLAYSTYLES.md --- Literature / Canadian Fish: play styles, rule dialects, and
engine specification.
\newblock Internal research synthesis with per-claim evidence tiers, 2026. Its Part~V
reports zero peer-reviewed prior work on Literature as a game-AI domain; several of its
secondary citations are flagged unverified and are not relied on here.

\bibitem{inertpaper}
A.~Hsieh.
\newblock The inert axis: when a style parameter is wired, swept, and never reached.
\newblock FishAI project; companion paper, published alongside this one, 2026.

\bibitem{containedpaper}
A.~Hsieh.
\newblock The contained book: an absorbing resource measured to be worth nothing.
\newblock FishAI project; companion paper, published alongside this one, 2026.

\bibitem{repo}
FishAI project.
\newblock FishAI repository: engine, bots, simulation laboratory, and results site.
\newblock README and audit record, 2026.

\bibitem{rules}
FishAI project.
\newblock RULES\_US54.md --- the ``US student'' 54-card variant.
\newblock Repository rules specification with derived consequences and test vectors, 2026.

\bibitem{styles}
FishAI project.
\newblock STYLES.md --- the play-style roster under \usff{}.
\newblock Repository document; \S3.1.1 the Hoarder inertness audit, \S6.1 the inert
declare-threshold axis, \S6.3 the contained-book turn-pass, \S6.4 matrix v2, 2026.

\bibitem{botlab}
FishAI project.
\newblock BOT\_LAB.md --- play-style spectrum: research, experimental design, and metrics.
\newblock Repository document; \S4.4 the verdict criteria, \S5 the experimental design,
\S10 threats to validity, 2026.

\bibitem{containment}
FishAI project.
\newblock CONTAINMENT.md --- the contained-book result.
\newblock Repository document; claims C1--C6 measured against the engine, 2026.

\bibitem{artifact}
FishAI project.
\newblock Tournament artifact \code{src/lab/data/style-results.v2.json}.
\newblock 36 cells $\times$ 4{,}300 duplicate pairs (309{,}600 games); records digest
\code{0c3cb524a3534d4a}; engine commit \code{1667a1d} (recorded \code{1667a1d-dirty});
generated 2026-08-22.

\end{thebibliography}

\end{document}
